\documentclass{article}
\usepackage[utf8]{inputenc}
\usepackage[spanish]{babel}
\usepackage{geometry}
\usepackage{tabularx}

\usepackage{longtable}
\usepackage{array}
\usepackage{colortbl} % Añadido para habilitar \rowcolor

\geometry{a4paper, margin=1in}

\title{PARQUE CIENTÍFICO Y TECNOLÓGICO UTPL \\
ISBEN Solutions – Plataforma Digital: ConectaMarket\\
Versión 1.0.0 \\
Manual de usuario }
\author{
Autores: \\
Líder de grupo: Byron Alvarez \\
Estudiante: Cody Cabrera \\
Estudiante: Byron Alvarez \\
Docente Tutor: Ing. René Elizalde
}
\date{}

\begin{document}

\maketitle

\section{Capítulo 1: Generalidades del Proyecto} 

\subsection{Introducción} 
El presente documento detalla el análisis y planificación para el desarrollo de una plataforma web de gestión de pedidos orientada a la comercialización de productos e insumos para micromercados. 
El sistema actúa como un puente tecnológico entre empresas mayoristas y tiendas locales (principalmente en zonas rurales), brindando oportunidades a vendedores independientes mediante un catálogo con inventario en tiempo real, incentivos por gamificación y un stack tecnológico moderno desarrollado de forma Full-Stack utilizando el framework web de Django con PostgreSQL para la persistencia de datos. 

\subsection{Problemática} 
Las empresas mayoristas y proveedores de micromercados en zonas rurales enfrentan una pérdida constante de ventas y clientes debido a un control deficiente y desfasado del inventario entre el punto de abastecimiento y la toma de pedidos en campo. 
Simultáneamente, la adopción de soluciones de compra directa (e-commerce) es inviable debido a la barrera tecnológica y desconfianza de los dueños de tiendas (población mayor a 60 años). 
Esto hace indispensable la figura de vendedores intermediarios, quienes, al no contar con información de stock en tiempo real ni incentivos dinámicos, generan incumplimientos en las entregas y estancamiento de productos. 

\subsection{Metodología o estrategia de desarrollo} 
Se empleará un enfoque ágil basado en Scrum, adaptado para el desarrollo iterativo e incremental del software. 
Esta metodología permitirá entregar módulos funcionales en ciclos cortos. La comunicación del equipo se gestionará mediante juntas de planificación de sprint, organizaciones diarias de sincronización (daily stand-ups) y juntas de revisión al finalizar cada iteración. 
Se utilizarán herramientas de control de versiones (Git) para el seguimiento de las tareas. 

\section{Capítulo 2: Planificación y Análisis} 

\subsection{Análisis} 
El proyecto se enfoca en resolver la brecha de información entre el inventario de la comercializadora y el vendedor en campo. 
Los interesados principales incluyen a las Comercializadoras (quienes pagan la suscripción y proveen el inventario), los Vendedores (usuarios que gestionan rutas, registran pedidos y ganan comisiones/puntos), y los dueños de Micromercados, quienes ahora tienen la capacidad de interactuar directamente con la plataforma mediante el rol de Tienda para auto-gestionar pedidos y llevar control de su inventario local. 
El alcance se centra en la intermediación, toma de pedidos y gestión de inventarios, dejando fuera la logística de entrega física. 

\subsection{Requisitos funcionales} 
\begin{table}[h]
\centering
\begin{tabular}{|l|p{12cm}|}
\hline
\textbf{ID} & \textbf{Descripción} \\ \hline
RF-01 & El sistema debe permitir el registro y autenticación de usuarios en la plataforma bajo los roles de Comercializadora, Vendedor y Tienda.  \\ \hline
RF-02 & La plataforma debe permitir a la Comercializadora gestionar su catálogo completo, creando productos, asignando categorías y controlando el inventario de bodegas.  \\ \hline
RF-03 & El sistema debe permitir a Vendedores y Tiendas visualizar el catálogo disponible y registrar pedidos verificando disponibilidad en tiempo real.  \\ \hline
RF-04 & El sistema debe calcular las comisiones de la plataforma y procesar la gamificación (asignación de puntos) al vendedor correspondiente de manera automática.  \\ \hline
RF-05 & El sistema debe gestionar una bandeja compartida de pedidos: los creados directamente por las Tiendas quedarán disponibles para ser asignados al primer Vendedor que los procese.  \\ \hline
RF-06 & El sistema debe actualizar el inventario propio de una Tienda (InventarioTienda) en cuanto un pedido pase a estado 'ENTREGADO'.  \\ \hline
\end{tabular}
\end{table}

\subsection{Requisitos no funcionales} 
\begin{table}[h]
\centering
\begin{tabular}{|l|p{12cm}|}
\hline
\textbf{ID} & \textbf{Descripción} \\ \hline
RNF-01 & Arquitectura: La plataforma debe estar construida de forma monolítica (Full-Stack) utilizando el sistema de plantillas de Django, delegando la lógica de negocio a los modelos (Fat Models, Skinny Views).  \\ \hline
RNF-02 & Base de Datos: Se utilizará PostgreSQL para garantizar la integridad, implementando bloqueo optimista (Optimistic Locking) para prevenir inconsistencias en el inventario.  \\ \hline
RNF-03 & Usabilidad: La interfaz renderizada por Django debe ser completamente responsiva (Mobile-First) con soporte nativo de CSS para su correcta visualización en campo.  \\ \hline
\end{tabular}
\end{table}

\subsection{Historias de usuario / Casos de Uso} 
\begin{itemize}
    \item HU-01 (Vendedor): Como vendedor, quiero ver el inventario actualizado en tiempo real para evitar vender productos que están fuera de stock y perder credibilidad con las tiendas. 
    \item HU-02 (Comercializadora): Como administrador de la comercializadora, quiero poder gestionar mi catálogo de productos y visualizar las liquidaciones pendientes para tener control financiero sobre las comisiones. 
    \item HU-03 (Vendedor): Como vendedor, quiero visualizar mis puntos acumulados en un dashboard para canjearlos y monitorear mis ganancias por comisiones. 
    \item HU-04 (Tienda): Como dueño de micromercado, quiero poder ingresar al sistema y generar pedidos directamente de los catálogos para abastecerme sin esperar la visita física de un vendedor.
    \item HU-05 (Tienda): Como dueño de micromercado, quiero consultar el estado de mis inventarios locales (InventarioTienda) y pedidos recientes para llevar un mejor control del negocio.
\end{itemize}

\subsection{Sprint} 
El desarrollo se dividirá en Sprints de 2 semanas: 
\begin{enumerate}
    \item Sprint 1: Diseño de base de datos en PostgreSQL, configuración de Django y creación del módulo de Autenticación/Roles (Vendedor, Comercio, Tienda). 
    \item Sprint 2: Desarrollo del CRUD de productos, categorías, inventarios e implementación de las interfaces de usuario con el motor de plantillas de Django. 
    \item Sprint 3: Módulo de toma de pedidos, lógica transaccional, cálculos automáticos de montos y manejo de bandeja compartida para tiendas. 
    \item Sprint 4: Implementación del motor de gamificación (Campañas Recompensa), liquidaciones y tableros (dashboards) finales para cada rol. 
\end{enumerate}

\subsection{Diagramas: Clases y BD (diseño lógico) - Preguntas y respuesta al diagrama} 
El diseño lógico en PostgreSQL contempla las siguientes entidades principales: Usuario, Vendedor, Comercializadora, Tienda, Categoria, Producto, Inventario, InventarioTienda, Pedido, DetallePedido, LiquidacionComercializadora, CampanaRecompensa, y TransaccionPuntos. 

\section{Diccionario de Datos (Atributos y Reglas de Negocio)}

A continuación se detalla el diccionario de datos derivado del diagrama de clases, especificando los tipos de datos a implementar en el ORM de Django y las reglas de negocio transaccionales para PostgreSQL.

\begin{longtable}{|>{\bfseries}p{3.2cm}|p{3.5cm}|p{2.8cm}|p{5.5cm}|}
\hline
\rowcolor[gray]{0.9} 
\textbf{Entidad} & \textbf{Atributo} & \textbf{Tipo (Django)} & \textbf{Descripción / Regla de Negocio} \\ \hline
\endfirsthead

\hline
\rowcolor[gray]{0.9} 
\textbf{Entidad} & \textbf{Atributo} & \textbf{Tipo (Django)} & \textbf{Descripción / Regla de Negocio} \\ \hline
\endhead

\hline \multicolumn{4}{|r|}{{Continúa en la siguiente página...}} \\ \hline
\endfoot

\hline
\endlastfoot

% --- USUARIO ---
Usuario & id & AutoField & Identificador único (PK). \\ \hline
Usuario & ruc & CharField(13) & Registro Único de Contribuyentes. Debe ser \texttt{unique=True} y tener exactamente 13 dígitos. \\ \hline
Usuario & nombres / apellidos & CharField(100) & Nombres y apellidos del usuario. \\ \hline
Usuario & email & EmailField & Correo electrónico. Debe ser \texttt{unique=True}. \\ \hline
Usuario & password & CharField(128) & Contraseña encriptada (Hash PBKDF2). \\ \hline
Usuario & estado\_cuenta & BooleanField & Indica si el usuario está activo (\texttt{default=True}). \\ \hline
Usuario & rol & CharField(20) & Opciones (Choices): 'VENDEDOR', 'COMERCIALIZADORA' o 'TIENDA'. \\ \hline

% --- VENDEDOR ---
Vendedor & usuario\_id & OneToOneField & PK y FK hacia Usuario. Relación 1 a 1. \\ \hline
Vendedor & zona\_asignada & CharField(100) & Región o ruta rural de cobertura. \\ \hline
Vendedor & vehiculo\_placa & CharField(10) & Placa del vehículo de transporte. \\ \hline
Vendedor & puntos\_acumulados & IntegerField & Total de puntos de gamificación. Regla: $\ge$ 0. \\ \hline
Vendedor & presupuesto\_ventas & DecimalField & Presupuesto asignado para compras al por mayor. Regla: $\ge$ 0.00. \\ \hline

% --- COMERCIALIZADORA ---
Comerciali\-zadora & usuario\_id & OneToOneField & PK y FK hacia Usuario. Relación 1 a 1. \\ \hline
Comerciali\-zadora & razon\_social & CharField(150) & Nombre legal de la empresa. \\ \hline
Comerciali\-zadora & nombre\_empresa & CharField(150) & Nombre comercial visible en el catálogo. \\ \hline
Comerciali\-zadora & direccion\_matriz & TextField & Ubicación principal del negocio. \\ \hline
Comerciali\-zadora & suscripcion\_activa & BooleanField & Define si sus productos se muestran en el catálogo. \\ \hline

% --- TIENDA ---
Tienda & id & AutoField & Identificador único (PK). \\ \hline
Tienda & usuario\_id & OneToOneField & PK y FK hacia Usuario. Relación 1 a 1 para propietarios de micromercados. \\ \hline
Tienda & nombre & CharField(100) & Nombre del micromercado. \\ \hline
Tienda & direccion & TextField & Dirección física en la zona rural. \\ \hline
Tienda & telefono & CharField(15) & Teléfono de contacto del micromercado. \\ \hline
Tienda & latitud / longitud & DecimalField & Coordenadas exactas para ubicar la tienda en el mapa. \\ \hline

% --- CATEGORIA ---
Categoria & id & AutoField & Identificador único (PK). \\ \hline
Categoria & nombre & CharField(50) & Nombre o clasificación general de familia de productos. \\ \hline

% --- PRODUCTO ---
Producto & id & AutoField & Identificador único (PK). \\ \hline
Producto & comercializadora\_id & ForeignKey & FK hacia Comercializadora. \\ \hline
Producto & codigo & CharField(50) & Código interno del producto. Debe ser \texttt{unique=True}. \\ \hline
Producto & nombre & CharField(150) & Nombre descriptivo del producto. \\ \hline
Producto & categoria\_id & OneToOneField & FK hacia Categoria. \\ \hline
Producto & precio & DecimalField & Costo base de venta al por mayor. \\ \hline

% --- INVENTARIO CENTRAL ---
Inventario & id & AutoField & Identificador único (PK). \\ \hline
Inventario & producto\_id & OneToOneField & FK hacia Producto (1 a 1 por bodega principal). \\ \hline
Inventario & almacen\_origen & CharField(100) & Nombre o código de la bodega. \\ \hline
Inventario & cantidad\_disp & Positive\-Integer\-Field & Stock en tiempo real. Regla: No puede ser menor a 0. \\ \hline
Inventario & fecha\_actualizacion & DateTimeField & Timestamp de actualización de stock. \\ \hline
Inventario & version & IntegerField & Control de concurrencia (Optimistic Locking). \\ \hline

% --- INVENTARIO TIENDA ---
Inventario\-Tienda & id & AutoField & Identificador único (PK). \\ \hline
Inventario\-Tienda & tienda\_id & ForeignKey & FK hacia Tienda a la que pertenece el stock. \\ \hline
Inventario\-Tienda & producto\_id & ForeignKey & FK hacia Producto adquirido. \\ \hline
Inventario\-Tienda & cantidad\_disp & Positive\-Integer\-Field & Unidades disponibles en la percha local de la tienda. \\ \hline

% --- PEDIDO ---
Pedido & id & AutoField & Identificador único (PK). \\ \hline
Pedido & vendedor\_id & ForeignKey & FK hacia Vendedor. Si es \texttt{null}, ingresa a bandeja compartida. \\ \hline
Pedido & tienda\_id & ForeignKey & FK hacia Tienda que recibirá los productos. \\ \hline
Pedido & fecha & DateTimeField & Fecha y hora exacta de creación (\texttt{auto\_now\_add=True}). \\ \hline
Pedido & estado & CharField(20) & Choices: 'PENDIENTE', 'CONFIRMADO', 'ENTREGADO', 'CANCELADO'. \\ \hline
Pedido & monto\_total\_tienda & DecimalField & Suma de subtotales. Total a cancelar. \\ \hline

% --- DETALLE PEDIDO ---
DetallePedido & id & AutoField & Identificador único (PK). \\ \hline
DetallePedido & pedido\_id & ForeignKey & FK hacia Pedido. Borrado en cascada. \\ \hline
DetallePedido & producto\_id & ForeignKey & FK hacia Producto. \\ \hline
DetallePedido & cantidad & Positive\-Integer\-Field & Unidades solicitadas. Regla: $>$ 0. \\ \hline
DetallePedido & precio\_unitario & DecimalField & Precio congelado al momento de la venta. \\ \hline
DetallePedido & subtotal & DecimalField & \texttt{cantidad * precio\_unitario}. \\ \hline

% --- LIQUIDACION COMERCIALIZADORA ---
Liquidacion\-Comerciali\-zadora & id & AutoField & Identificador único (PK). \\ \hline
Liquidacion\-Comerciali\-zadora & pedido\_id & OneToOneField & FK hacia Pedido. Generado desde la clase modelo de Pedido. \\ \hline
Liquidacion\-Comerciali\-zadora & monto\_comision & DecimalField & Fee de la plataforma (10\% del total por defecto). \\ \hline
Liquidacion\-Comerciali\-zadora & monto\_cobrar & DecimalField & Lo que recibe la comercializadora (Total - Comisión). \\ \hline
Liquidacion\-Comerciali\-zadora & fecha\_liquidacion & DateTimeField & Fecha en que se procesó el corte contable. \\ \hline
Liquidacion\-Comerciali\-zadora & estado\_pago & CharField(20) & Choices: 'PENDIENTE', 'PAGADO\_COMERCIALIZADORA'. \\ \hline

% --- CAMPANA RECOMPENSA ---
Campana\-Recompensa & id & AutoField & Identificador único (PK). \\ \hline
Campana\-Recompensa & producto\_id & ForeignKey & FK hacia Producto. \\ \hline
Campana\-Recompensa & nombre\_campana & CharField(100) & Identificador de la estrategia (Ej: 'Impulso Lácteos'). \\ \hline
Campana\-Recompensa & factor\_puntos & Positive\-Integer\-Field & Cantidad de puntos extra que da este producto al venderse. \\ \hline
Campana\-Recompensa & fecha\_inicio & DateTimeField & Inicio de vigencia de la campaña. \\ \hline
Campana\-Recompensa & fecha\_fin & DateTimeField & Fin de vigencia de la campaña. \\ \hline
Campana\-Recompensa & estado & BooleanField & Activa/Inactiva. \\ \hline

% --- TRANSACCION PUNTOS ---
Transaccion\-Puntos & id & AutoField & Identificador único (PK). \\ \hline
Transaccion\-Puntos & vendedor\_id & ForeignKey & FK hacia Vendedor. \\ \hline
Transaccion\-Puntos & pedido\_id & ForeignKey & FK hacia Pedido. Puede ser \texttt{null=True} si es un canje. \\ \hline
Transaccion\-Puntos & tipo\_transaccion & CharField(15) & Choices: 'INGRESO' (Venta), 'EGRESO' (Canje de apuestas). \\ \hline
Transaccion\-Puntos & puntos\_ganados & IntegerField & Valor absoluto de la transacción. \\ \hline
Transaccion\-Puntos & fecha & DateTimeField & Timestamp de la transacción. \\ \hline

\end{longtable}

A continuación, se plantean preguntas críticas para validar la robustez del modelo de clases y los requerimientos funcionales actuales: 

\textbf{Pregunta 1 (Concurrencia):} ¿Cómo resuelve el modelo de clases la concurrencia si dos vendedores intentan registrar un pedido para el mismo producto simultáneamente, y el stock disponible es insuficiente para ambos?  \\
\textbf{Respuesta:} El modelo en Django centraliza la lógica dentro de una transacción atómica mediante un control de bloqueo optimista utilizando el campo \texttt{version} en la entidad Inventario. Antes de guardar el DetallePedido, la clase del modelo evalúa la disponibilidad y detiene la petición levantando una alerta \texttt{StockInsuficienteError} si no existe stock real. 

\textbf{Pregunta 2 (Acoplamiento de Pagos):} Si el modelo de negocio cobra una comisión del 10\% por venta, ¿dónde se registra esta obligación para no afectar el monto bruto que la tienda debe pagar a la comercializadora?  \\
\textbf{Respuesta:} La clase Pedido posee la función interna \texttt{generar\_liquidacion()} para automatizar esto. Genera el registro \texttt{LiquidacionComercializadora} de forma desacoplada para mantener la trazabilidad de los pagos por parte de los administradores. 

\textbf{Pregunta 3 (Toma de Auto-Pedidos):} ¿Cómo se administra un pedido generado directamente por el rol de Tienda? \\
\textbf{Respuesta:} Cuando un propietario de tienda genera su pedido, este se crea dejando el campo \texttt{vendedor} nulo. A través de consultas centralizadas, este pedido cae en una "bandeja compartida" visualizada por todos los Vendedores; el primero que procese o administre el pedido asume el rol del \texttt{vendedor\_id} y de las comisiones.

\subsection{Diagrama de componentes} 
El sistema se compone de: 
\begin{itemize}
    \item Componente de Interfaz de Usuario (Django Templates): Sistema de renderizado de pantallas con HTML/CSS mediante vistas y formularios acoplados directamente al backend. 
    \item Componente Core de Negocios (Django Models/Views): Módulo central de toda la aplicación que unifica la lógica relacional con base en roles y operaciones CRUD (Create, Read, Update, Delete). 
    \item Componente de Persistencia (PostgreSQL): Motor de bases de datos para el almacenamiento seguro de relaciones transaccionales, control de stock y usuarios. 
\end{itemize}

\subsection{Arquitectura lógica y física (propuesta)} 
\textbf{Arquitectura Lógica:} Modelo Cliente-Servidor (Aplicación web monolítica basada en renderizado del lado del servidor - SSR - gestionada íntegramente por Django y su patrón MVT).  \\
\textbf{Arquitectura Física:} 
\begin{itemize}
    \item Capa Cliente: Navegadores web en dispositivos móviles y computadoras de escritorio visualizando HTML enriquecido. 
    \item Capa de Presentación / Proxy: Servidor Nginx que actúa como proxy inverso y sirve los recursos estáticos del sistema de plantillas. 
    \item Capa de Aplicación: Contenedores Docker ejecutando Gunicorn para levantar y procesar la aplicación principal de Django de forma estable. 
    \item Capa de Datos: Servidor de base de datos relacional PostgreSQL con protección ante concurrencia transaccional. 
\end{itemize}

\subsection{Prototipo de pantallas no funcional}

Con el propósito de validar la experiencia de usuario y visualizar el alcance funcional de la plataforma antes de su implementación definitiva, se desarrolló un conjunto de prototipos no funcionales. Estas interfaces permiten representar los principales módulos del sistema y los flujos de interacción de los distintos tipos de usuarios.

\begin{itemize}
	
	\item \textbf{Página de Inicio:} Pantalla principal de presentación de la plataforma, donde se describen los beneficios para comercializadoras, vendedores y micromercados, además de proporcionar acceso a las opciones de inicio de sesión y registro.
	
	\item \textbf{Pantalla de Inicio de Sesión:} Permite la autenticación segura de los usuarios registrados (Vendedores, Comercializadoras, Tiendas).
	
	\item \textbf{Pantalla de Registro:} Facilita el registro de nuevos usuarios bajo los tres roles habilitados, solicitando la información necesaria para su creación inmediata en base de datos.
	
	\item \textbf{Dashboard (Panel de Control Principal):} Vistas específicas para cada rol. El Comercio ve sus productos; el Vendedor analiza sus pedidos y comisiones; y la Tienda consulta su inventario local y estado de reabastecimiento.
	
	\item \textbf{Catálogo de Productos:} Presenta los productos disponibles agrupados por categoría comercial con visibilidad de precios y reglas de inventario.
	
	\item \textbf{Gestión de Inventario (B2B y B2C):} Permite al comercio mayorista abastecer su bodega principal, mientras a la tienda local le permite controlar su stock en percha.
	
	\item \textbf{Registro de Pedido:} Interfaz unificada usada tanto por Vendedores de ruta como por dueños de Tienda para generar pedidos calculando subtotales automáticamente.
	
	\item \textbf{Seguimiento y Transición de Pedidos:} Módulo destinado a consultar y avanzar el estado de pedidos. Permite inyectar inventario automáticamente hacia las tiendas una vez que son 'ENTREGADOS'.
	
	\item \textbf{Panel de Comisiones y Recompensas:} Tablero para que los vendedores consulten el efecto financiero de cada venta generada.
	
\end{itemize}

Cabe destacar que el sistema actualmente ya no es un prototipo y se encuentra en etapa funcional. La estructura visual renderizada hereda la maquetación definida en el sistema de plantillas y CSS base provistos en los componentes de Django.

\end{document}